\chapter{Future Work}
\label{sec:future-work}

This chapter describes three directions that extend \project{}'s architecture beyond the scope of this thesis: integrating a high-level query language, automating layout and property mode selection through learned optimization, and serving graph machine learning workloads directly from the storage layer.

%%%%%%%%%%%%%%%%%%%%%%%%%%%%%%%%%%%%%%%%%%%%%%%%%%%%%%%%%%%%%%%%%%%%%%%%%%%%%%%%
\section{High-Level Query Language Integration}
\label{sec:future:query-lang}

\project{} currently exposes only a C++ cursor API.
All queries in the evaluation chapters (\autoref{sec:eval:inmem}--\autoref{sec:eval:dynamic}) are written as imperative programs that open cursors, position them, and iterate.
This low-level interface gives fine-grained control over the access path, but it places the burden of query planning entirely on the programmer.

A declarative query language would let users express graph queries without specifying how to traverse the underlying tables.
The natural target is GQL, the ISO standard for property graph queries~\cite{francis2023gql_standard}, or openCypher~\cite{openCypher}, which is already supported by Neo4j, Memgraph, and several other systems.
Compilation would map declarative pattern matching and path expressions down to \project{}'s cursor-level operations: vertex iterators for full scans, neighborhood iterators for adjacency traversal, and property cursors for predicate evaluation.

The main challenge is query planning and optimization.
Traditional graph databases maintain a catalog of statistics (label cardinalities, degree histograms, index selectivity estimates) and use these to choose join orderings and access paths.
\project{} does not maintain such a catalog; \wt{} provides low-level statistics (table size, cache residency), but not the graph-level statistics a query planner needs.
Building a lightweight catalog that tracks per-label vertex counts, degree distributions, and property value histograms is a prerequisite for cost-based query planning.

A second question is whether a language model can serve as a stopgap query planner.
Recent work on LLM-based text-to-SQL and text-to-Cypher translation~\cite{neo4j_neodash} demonstrates that language models can generate syntactically correct graph queries from natural language.
Extending this approach to \project{} would require the model to select between AdjList and EdgeKey access paths, choose between embedded and columnar property reads, and decide on scan direction (in-neighbors vs.\ out-neighbors) based on the workload context.
Whether an LLM-assisted approach can produce plans competitive with a conventional cost-based optimizer is an open question, but it could inform the design of a future query planner by identifying which plan choices matter most in practice.


%%%%%%%%%%%%%%%%%%%%%%%%%%%%%%%%%%%%%%%%%%%%%%%%%%%%%%%%%%%%%%%%%%%%%%%%%%%%%%%%
\section{Learned Indexes and Adaptive Layout Selection}
\label{sec:future:learned}

A central claim of this thesis is that data layout is the primary performance determinant for graph analytics (\autoref{chap:design_space}).
The evaluation chapters confirm this: AdjList outperforms EdgeKey on read-heavy analytics (\autoref{sec:eval:inmem}), while EdgeKey provides faster transactional mutation (\autoref{sec:eval:dynamic}); columnar property mode outperforms embedded mode on selective column access (\autoref{sec:architecture:property}).
However, the user must currently select the structural representation and property mode at graph creation time, and this choice is fixed for the lifetime of the graph.

The system could instead observe the query workload and data distribution and adapt its physical layout automatically.
There are four concrete decision points amenable to learned optimization:

\begin{enumerate}
\item \textbf{AdjList vs.\ EdgeKey selection.}
The read/write ratio and the degree distribution of each label determine which structural representation is more efficient.
A workload monitor could track the ratio of traversal operations to edge insertions and deletions per label and recommend switching representations when the observed ratio crosses a threshold.
The checkpoint-based migration path described in Section~\ref{sec:future:learned:migration} would enable such a switch without downtime.

\item \textbf{Embedded vs.\ columnar property mode.}
The choice depends on which properties are co-accessed by queries.
If most queries access all properties of a vertex (e.g., full profile lookups), embedded mode avoids the overhead of opening multiple column group cursors.
If queries consistently access a single property column (e.g., scanning creation dates), columnar mode avoids reading unneeded columns.
A per-label workload profile that tracks which property columns are accessed together would identify when to switch modes.

\item \textbf{Column group composition.}
Within columnar mode, the assignment of properties to column groups determines I/O isolation.
Properties that are frequently co-accessed should share a column group (one B-tree read retrieves both), while properties accessed independently should be in separate column groups.
This is the graph-property analogue of the column-group selection problem in column stores~\cite{copeland1985decomposition}, and query co-access frequency matrices provide the input signal for optimization.

\item \textbf{Cursor prefetching hints.}
\project{}'s topology cursors iterate over adjacency structures stored in \wt{} B-trees.
If the system can predict the next traversal step (e.g., a BFS will access all neighbors of the current frontier), it can issue prefetch hints to \wt{}, bringing pages into the cache before they are needed.
Learned traversal pattern models could generate such hints, reducing cache miss stalls during analytics.
\end{enumerate}

\subsection{Checkpoint-Triggered Layout Migration}
\label{sec:future:learned:migration}

Adaptive layout selection requires a mechanism for migrating data between representations without stopping the system.
\project{} can support this because its \code{GraphBase} API abstracts over both AdjList and EdgeKey: application code that uses topology cursors does not depend on the underlying representation.

A migration procedure tied to \wt{}'s checkpoint mechanism would work as follows.
At checkpoint time, \wt{} writes a consistent snapshot to disk.
A background process reads from the source representation in the checkpoint snapshot and bulk-loads the data into the target representation, using the existing bulk loader infrastructure.
On completion, the engine swaps the active graph handle from the source representation to the target.
New writes resume in the target representation; read queries that were in flight against the old representation drain against the snapshot.

The cost model for such a migration is an open question.
The research contribution would be to identify the breakeven point: for a given read-to-write ratio, the analytics speedup from switching to AdjList must exceed the migration cost.
Validating this cost model against LDBC SNB workloads that mix Interactive (write-heavy) and BI (analytics-heavy) phases would quantify when migration pays for itself.

A complication is that migration introduces a lag between the source and target representations.
Writes that arrive after the migration checkpoint exist only in the source; queries against the target must merge results from both representations until the next migration cycle.
This dual-format read problem is analogous to the stable-plus-durable view merging in log-structured systems like GraphOne~\cite{graphone2019kumar}.

Related work on adaptive data structures includes the learned index structures of Kraska et al.~\cite{kraska2018case}, SageDB~\cite{kraska2019sagedb}, and Bao~\cite{marcus2021bao} (learned query optimization for relational databases).
Extending these ideas to graph-specific layout decisions, where the choice space includes topology format, property mode, and column group assignment, is an open area with no prior work.


%%%%%%%%%%%%%%%%%%%%%%%%%%%%%%%%%%%%%%%%%%%%%%%%%%%%%%%%%%%%%%%%%%%%%%%%%%%%%%%%
\section{Graph Storage for Machine Learning Workloads}
\label{sec:future:gnn}

Graph Neural Networks (GNNs) have become the dominant approach for machine learning on graph-structured data, with applications in fraud detection, recommendation systems, drug discovery, and knowledge graph completion~\cite{hamilton2017graphsage, velickovic2018gat, kipf2017gcn}.
GNN training and inference impose specific demands on graph storage that current practice handles with ad hoc solutions: DGL~\cite{wang2019dgl} and PyG~\cite{fey2019pyg} maintain their own in-memory graph representations, while production systems add separate graph sampling services in front of the primary store.
\project{}'s architecture addresses these requirements without such auxiliary infrastructure.

\subsection{Training: Neighborhood Sampling on Large Graphs}
\label{sec:future:gnn:training}

GNN training proceeds via mini-batch neighbor sampling.
For each mini-batch, the system selects a batch of target nodes, samples $k$ neighbors per node for 2--3 hops to construct a computation subgraph, loads the feature vectors (node properties) for every node in the subgraph, and runs the neural network forward and backward passes.
The access pattern is highly random: random target selection followed by pointer-chasing through the adjacency structure, then a batch of random point lookups for feature vectors.
This is closer to OLTP than OLAP and favors B-tree point lookups over sequential CSR scans.

Current practice requires loading the entire graph into memory or using framework-specific out-of-core storage.
Disk-based GNN training is an active area of research.
HUGE~\cite{yang2022huge}, GNNLab~\cite{yang2022gnnlab}, MariusGNN~\cite{mohoney2023mariusgnn}, and BGL~\cite{liu2023bgl} address neighbor sampling and feature loading from disk, but all assume a static graph snapshot loaded before training begins.

\project{} can serve this workload directly.
For graphs exceeding main memory, \project{}'s out-of-core capability provides snapshot-consistent neighborhood iteration through \wt{}'s buffer pool, without requiring the entire graph to reside in memory.
Columnar property mode fits the training access pattern: loading a single feature column (e.g., node embeddings or edge weights) across many vertices in a mini-batch requires only a sequential scan of the corresponding column group B-tree, touching none of the unneeded property columns.

Snapshot isolation adds a second benefit.
Each mini-batch can read from the same \wt{} snapshot timestamp, giving reproducible training semantics even as the graph is being concurrently updated.
This \emph{transactional GNN training}, consistent snapshots per mini-batch on a mutating graph, is not possible with CSR-based systems (HUGE, GNNLab) that lack MVCC, nor with in-memory frameworks (PyG, DGL) that lack persistence.
No existing disk-based GNN system combines durable out-of-core storage with MVCC-consistent reads.


\subsection{Inference: Low-Latency Neighborhood Lookups on Live Graphs}
\label{sec:future:gnn:inference}

Real-time GNN inference, for example classifying a new transaction as fraudulent or generating a recommendation for an active user, requires low-latency neighborhood lookups on a graph that is simultaneously being mutated.
This is the \ac{HGTAP} use case: transactional writes (new edges, updated properties) proceed concurrently while inference queries read consistent snapshots via MVCC.

\project{}'s property cursors provide the feature vectors needed for inference without a separate feature store.
An inference query opens a read transaction at the current timestamp, samples the target node's neighborhood using topology cursors, loads the required feature columns via property cursors, and produces an embedding.
The entire pipeline runs against a single consistent snapshot of the graph, with no coordination with concurrent writers beyond \wt{}'s standard MVCC protocol.


\subsection{Open Problems}
\label{sec:future:gnn:open}

Serving GNN workloads from \project{} raises several concrete engineering and research problems.

\paragraph{Language bridge.}
\project{} has a C++ API; ML frameworks (PyTorch, JAX) are Python-based.
Bridging the two requires a Python binding layer.
The PyTorch Geometric (PyG) framework defines two abstract interfaces, \texttt{FeatureStore} for node properties as tensors and \texttt{GraphStore} for edge indices in COO format, that enable a graph database to serve as a remote backend for GNN training~\cite{fey2023pyg_remote_backend}.
Implementing these interfaces for \project{} would involve a pybind11 layer that exposes \wt{} cursor scans to Python via NumPy buffers: scanning a property column fills a pre-allocated NumPy array in place, and \texttt{torch.from\_numpy()} shares the memory without copying.
Edge extraction would scan the adjacency column with a \wt{} cursor, batch results into a pre-allocated array, and reshape to the $[2, |E|]$ COO tensor format that PyG expects.

\paragraph{Batched sampling.}
GNN training requires batches of multi-hop neighborhoods, not single traversals.
\project{}'s current iterator API processes one traversal at a time: the caller opens a neighborhood cursor, iterates over neighbors, and closes the cursor.
A batched sampling primitive that accepts a set of source nodes and a fan-out specification ($k$ neighbors per hop for $h$ hops) and returns the entire computation subgraph in one call would reduce per-cursor overhead and enable I/O coalescing across the batch.

\paragraph{Feature materialization.}
Some GNN architectures require pre-computed aggregate features (e.g., mean neighbor embeddings, degree-weighted sums) as inputs.
Materializing these derived features as additional property columns in \project{}'s columnar store would avoid redundant computation during training.
The materialized columns would function as graph-aware materialized views: invalidated when the underlying topology or source properties change, and refreshed either eagerly (on write) or lazily (on next read transaction).

\paragraph{Embedding storage and vector similarity.}
After training, GNN models produce dense embedding vectors for each node (typically 128--768 dimensions, 512\,B--3\,KB per node).
Storing embeddings as a \wt{} property column works for point lookups by node ID, which covers the training and inference access pattern.
However, similarity queries (``find the $K$ nodes most similar to this one'') require approximate nearest neighbor (ANN) search, which cannot be expressed as a B-tree range scan.
An in-process side-car HNSW index~\cite{malkov2018hnsw}, rebuilt from a \wt{} snapshot and atomically swapped into service, would provide sub-millisecond ANN queries with snapshot-consistent semantics.
No other graph database currently supports this.
